\chapter{Grammar Was Not Invented by Teachers}
\section{A Red Pen}
First pick up a red pen.

Not the sort you buy at the stationery shop, but the one in your teacher's drawer: tooth marks on the cap, clear tape around the barrel, ink that flows fiercely and lands on paper like a little pool of blood. It has one job: to pass judgment.

Consider a few recent cases. A student writes, ``Zhè bù diànyǐng wǒ yǒu kànguo''---``This film, I have seen,'' using yǒu, ``have,'' before the experiential verb. The red pen circles yǒu and writes: faulty sentence; delete ``have.'' Another writes, ``It's getting late; nǐ zǒu xiān ba,'' literally ``you leave first,'' with xiān, ``first,'' after the verb. The pen circles zǒu xiān and changes it to xiān zǒu, ``first leave.'' A third writes, ``He is taller than me by one head,'' yī ge tóu. This time the pen hesitates. Surely that is fine? But it circles the phrase anyway: the workbook's standard answer says yī tóu, ``one head,'' without the classifier ge.

Notice something strange. A Chinese speaker understands all three convicted sentences perfectly, without looking up a word. Their meanings could hardly be clearer. Moreover, thousands upon thousands of people actually say them. In Guangdong teahouses, ``you leave first'' in that order is spoken tens of thousands of times a day. Nobody mistakes it for ``lean your body forward when walking.''

So here is the question: where is the ``illness'' in a sentence everyone understands and people say daily? Who issued its diagnosis, and on what basis?

Pursue the question, and even the doctor may struggle to explain the pathology. Ask a teacher why ``I leave first,'' with first after the verb, is wrong, and the common answer is: ``Grammar requires it.'' Ask, ``Who required it?'' and the teacher may pause before saying, ``That's how everyone speaks.'' But a hundred million people in Guangdong could raise their hands: our ``everyone'' does not speak that way.

Behind this red pen stands an assumption accepted since childhood and rarely tested: grammar is a rulebook kept by experts and teachers, while our mouths must obey. In this chapter we will lay that assumption on the table and take it apart. You may find that real grammar does not live in workbooks at all. It lives in your brain, in a hundred million Guangdong speakers' teahouses, in three- and four-year-olds' mouths, operating more precisely than any grammar book can describe.

\section{Friday Afternoon in Chinese Class}
Now come into a scene.

Time: a Friday in September, after the second afternoon lesson. Place: a middle-school classroom in a southern city. The ceiling fan turns too slowly to stir the room's heat; cicadas call in waves outside. The duty student's name is written at the blackboard's upper right, and chalk dust drifts through slanting sunlight.

The lesson covers faulty sentences. Teacher Zhou is in her forties, quick and steady with chalk. She writes a question from the practice sheet:

``It's getting late; nǐ zǒu xiān ba''---``you leave first,'' with first following leave.

``Who can correct it?'' she asks.

The class representative raises a hand. ``Change it to nǐ xiān zǒu ba, `you first leave.' Xiān is an adverbial. It must precede the verb zǒu, not follow it.''

``Correct.'' Teacher Zhou nods. ``Adverbials go first. That is the standard word order.''

In the middle of the room, a hand rises hesitantly, stops halfway, then shrinks back a little. Teacher Zhou sees it. It belongs to Ah Chan, newly transferred from Guangzhou three weeks into term. His Mandarin is slow, each word seemingly rehearsed silently before he speaks.

``Teacher,'' he says, standing, ``why can't xiān go afterward?''

``Because standard modern Chinese puts the adverbial before the verb. Zǒu xiān does not conform to grammar.''

``But,'' Ah Chan says quietly, in a room so still everyone hears, ``where I'm from, everyone says `you leave first' that way. At tea we say `you drink first'; offering a seat, `you sit first.' Nobody thinks it's wrong. Are they\ldots{} all wrong?''

The fan creaks. Teacher Zhou pauses for two seconds, chalk in hand.

``In exams,'' she says, ``follow the standard.''

The bell rings. Chair legs scrape across cement; classmates stream toward the door. Ah Chan stays seated. He is not angry at being rebuffed---Teacher Zhou was quite gentle---but he noticed that she did not answer his question. ``Follow the standard in exams'' answers ``What should I write on the test?'' He asked something else: how can rules followed by a hundred million people become ``ungrammatical''? What, exactly, is operating in those hundred million brains when they speak?

Remember that question hanging in the air. Most of the rest of this chapter answers it.

\section{Who Made These Rules?}
First lay out the conflict, without rushing to an answer.

Teacher Zhou has a case: Chinese has standard grammar, adverbials precede verbs, and this is printed plainly in syllabuses and grammar books. Without standards, how do we grade exams, proofread newspapers, or understand people from elsewhere? Standards are not somebody's whim; they give hundreds of millions of speakers a common well to draw from.

Ah Chan has a case too: Cantonese ``you leave first'' is not merely understood by everyone, but used consistently. Postposed ``first'' means ``do this before the other things, which can wait,'' just like Mandarin's ``you first leave.'' This is plainly a functioning set of rules. Why call its sentences ``faulty''?

Behind both positions lies a real question:

\textbf{Where do grammatical rules come from?}

Roughly, there are just three candidates.

First: authority makes them. Scholars, academies, and education departments formulate grammar, put it in books, and everyone obeys, just as traffic authorities make traffic rules. The red pen enforces this system.

Second: they grow naturally. Language ability is equipment built into the brain; grammar is how it operates, like the beating of a heart or image formation in the eye. Nobody needs to decree it.

Third: they emerge through use. Grammar is habit deposited by billions of utterances, like a path worn through grass: no blueprint, but enough walkers create a route, and it can bend.

Notice a fact that troubles the authority camp: people do not learn to speak by first studying grammar books. Three- and four-year-olds, without a single grammar lesson, already produce complex sentences they have never heard. Their ``mistakes'' are even more revealing. An English-learning child says ``I goed to school,'' though no adult says goed; the child invented it. A Chinese-learning child overextends classifiers, calling a horse yī zhī mǎ, using zhī, a classifier for many animals, instead of the conventional horse classifier. These errors are crucial evidence: the child is not memorizing sentences but \textbf{extracting rules}---``add -ed for the past,'' ``use zhī for animals''---then applying them to every word, including exceptions. Memorization cannot produce these errors; rule-making produces such ``beautiful mistakes.''

Consider a more basic problem. If authority creates grammar, what did humans speak during the hundreds of thousands of years before writing, schools, or the profession of ``expert''? Archaeologists' earliest writing is only a little over five thousand years old, while speaking humans go back at least a hundred thousand years. Throughout that immense span, every sentence followed precise grammar. Who promulgated it?

Conversely, if grammar is wholly inborn or emerges from use, why does the red pen exist? Why does every society work so hard to teach children to speak and write ``correctly''?

Before reading on, choose your sympathies: Teacher Zhou or Ah Chan?

\section{The Red Pen's Case and the Ear's Case}
Let us state both positions fully, then see how other places handle the same difficulty.

\textbf{Position one: language needs standards, and the red pen does real work.}

Give Teacher Zhou's side its full case. This deserves care: ``standards'' can sound like an insult during your rebellious years, which is unfair to them.

The standardizers have at least three reasons. First, common standards are infrastructure for communication on a large scale. China has more than a billion people and hundreds or thousands of dialects. Without a shared standard language, someone from Heilongjiang and someone from Guangxi could communicate only through gestures. Promoting Mandarin is like standardizing rail gauges and electrical sockets: not to eliminate other gauges, but to let trains run nationwide. Second, standards give learning a reliable line. Learning characters, essay writing, or translation requires clear answers: ``this is correct; that is not.'' If everything becomes ``write however you like,'' education has nowhere to start and exams cannot be marked. Third, ambiguity can be dangerous. Statutes, medicine instructions, contracts, and flight manuals need sentences with a single interpretation, backed by stable standards.

This is no weak case. Many societies take it seriously. France established the Académie française in 1635; among its forty ``immortals''' central tasks are compiling dictionaries, setting standards, and ruling on ``correct French.'' It still meets and votes today. From the 1950s, China systematically promoted Mandarin, simplified characters, and developed Hanyu Pinyin, transforming literacy and cross-regional communication. Look further back: grammarians held extraordinarily high status in India more than two millennia ago. Panini's Ashtadhyayi, the ``Eight Chapters,'' packed Sanskrit's structure into a precise ``rule machine'' of nearly four thousand rules, whose rigor still amazes computer scientists. Notice that Panini was actually \textbf{recording} a living language. Later generations treated his book as legislation, ``freezing'' Sanskrit into its rules as an unchanging classical language. This reminds us that once description becomes sacred decree, a living language becomes a specimen.

\textbf{Position two: grammar already grows in users' mouths; books only photograph it.}

The other side's case is equally complete. Linguists, who study how language works, set themselves this rule: study language as biologists study butterflies. A biologist's job is not to teach butterflies ``correct flying posture,'' but to record how they actually fly. If a hundred million people say ``you leave first,'' the linguist concludes not ``a hundred million people are wrong,'' but ``our grammar book omitted a rule.''

This is not contrariness; there is hard evidence. Consider a frequently mocked ``error'': multiple negation in English, such as ``I can't get no nothing,'' found in a number of dialects and ridiculed as ``uneducated speech.'' Linguists recording these varieties find their negation meshes like gears: when to use it and what force it expresses are consistent throughout the community. It is a \textbf{rule}, simply not the textbook's rule. More striking still, standard French \textbf{requires} two negative words, ne\ldots{}pas; omitting one is ungrammatical. ``I don't know'' must be Je ne sais pas, with both negative components. The same ``double-negative'' apparatus is canonical in French and a laughingstock in English dialects. The difference is not in the rules themselves, but in whose rules enter textbooks.

Dialects make the point even clearer. Cantonese is not ``badly pronounced Mandarin,'' but a system of its own. Mandarin uses le and guo to mark aspects of actions; Cantonese distributes the work among particles including zo2, gan2, and gwo3. Ngo5 sik6 zo2 faan6, ``I have eaten,'' differs clearly from ngo5 sik6 gan2 faan6, ``I am eating.'' Some distinctions cannot be made the same way in Mandarin. Southern Min, Shanghainese, Hakka: dig into any of them and you find a structurally complete underground palace. A ``faulty sentence'' spoken in dialect is often entirely legitimate within its speaker's grammatical system.

Thus the disagreement is not whether to have rules, but \textbf{what status rules have}. One side treats them as laws requiring enactment and enforcement; the other as laws of physics. Apples need no enforcement to fall; gravitation describes what already happens. Grammarians name the two approaches: \textbf{prescriptive grammar} tells you ``how you should speak''; \textbf{descriptive grammar} records ``how you actually speak.'' The red pen holds the former, the ear the latter. Our title, ``Grammar Was Not Invented by Teachers,'' does not mean standards are useless. It means grammar was already operating long before standards appeared.

\section{Thinking Bridge: Three Glues for Sentences}
Here is a tool to take away. First correct an intuition: we tend to imagine words lining up in sentences like people queuing for milk tea. They do not. Words enter sentences in \textbf{groups}, like building blocks: several words form a small unit, then several units a larger one.

Ambiguity is the evidence. Consider the Chinese sequence yǎosǐ le lièrén de gǒu, roughly ``bit-to-death hunter's dog.'' It has two meanings: the hunter associated with the dog is the victim (the dog killed the hunter by biting), or that hunter is the killer (the hunter bit the dog to death). The same sequence of characters yields completely different meanings. If words simply queued, one sequence should have one meaning. Ambiguity shows that the same characters can build different ``block structures.'' Grasp this, and you gain X-ray vision for sentences.

Now for the tool itself. Across thousands of languages, the main ways of building words into sentences amount to \textbf{three kinds of glue}.

\textbf{First glue: word order.} Changing who goes first changes meaning. Chinese and English rely heavily on it: ``cat chases dog'' and ``dog chases cat'' contain identical words; order alone identifies the pursuer. Japanese places the verb last---``I rice eat''---a different order from Chinese, yet equally systematic. The price of inexpensive word order is rigidity: positions cannot move freely.

\textbf{Second glue: inflection.} Words grow their own parts: changing endings carry identity tags. Latin is a major user. In Petrus amat Paulam, ``Peter loves Paula,'' the -am ending on Paulam marks the object. Because the tag lives on the word, rearranging the words leaves the meaning intact: Paulam amat Petrus still means Peter loves Paula. English retains traces: he and him, I and me, are fossils of subject and object tags. Russian, German, and Arabic still display an abundant array of this machinery.

\textbf{Third glue: function words.} If words carry no tags of their own, little specialist words serve as signposts. Chinese bǎ, bèi, le, zhe, guo, and de are such signs. In wǒ bǎ bēizi dǎ le, ``I hit the cup,'' bǎ announces in advance that the cup receives the action. In tā bèi dǎbài le, ``he was defeated,'' bèi reverses the action's direction. These words refer to no thing---you cannot draw a portrait of le---yet remove them and the sentence falls apart. Japanese wa, ga, and o are the same kind of apparatus.

Every language mixes the three glues in its own proportions: Chinese favors order and function words, Latin heavy inflection, English lies between. None is more ``advanced.'' They put the same information in different containers.

With these glues, you can perform ``linguistic intuition experiments.'' Your brain contains a grammar book you have never read. Three little operations make it visible:

\begin{itemize}
\item \textbf{Substitution:} replace a word. Does the sentence still work? Wǒ chī fàn, ``I eat rice,'' can become wǒ chī fàncài, ``I eat food,'' but not wǒ fàn zhuōzi, ``I rice table.'' Interchangeable words belong to a class. You have just felt the boundary between ``noun'' and ``verb'' yourself, without a definition.
\item \textbf{Movement:} move one building block. Does meaning change? Turn ``cat chases dog'' into ``dog chases cat,'' and the meaning reverses. You have tested what word order controls in Chinese.
\item \textbf{Adding or removing parts:} add le or substitute bèi. How do meaning and the sense of time change? How does tā lái le, ``he has come,'' differ from tā láizhe, roughly ``he came'' with retrospective coloring? You have tested what function words control.
\end{itemize}
Next time you hear an ``ungrammatical'' sentence, do not rush to judgment. Use these three operations: has it really fallen apart, or merely used \textbf{another set} of glues?

\section{A Sentence from Xunzi and a Mountain Temple}
Now read a short primary text. More than 2,300 years ago, someone was already asking who makes the rules.

Xunzi, a thinker of the late Warring States period, wrote ``Rectification of Names,'' discussing where names and expressions come from. It contains this sentence:

\begin{quote}
Names have no inherent appropriateness: they are assigned by agreement. When agreement becomes established custom, they are called appropriate; what departs from agreement is called inappropriate. (Xunzi, ``Rectification of Names'')
\end{quote}
In essence: nothing is naturally the ``right'' name for a thing. People agree on a name, and when agreement becomes custom, it is appropriate. What violates that agreement is inappropriate.

Take this apart and weigh it. Xunzi is discussing names---why a dog is called a dog---but the principle extends to grammar. Why must ``first'' precede the verb? No heavenly principle, physical law, or invention patent requires it. There is only a silent collective vote lasting countless generations. Each generation is born into the polling place, learns what it hears, then votes for the next. Rules draw their strength not from a legislator but from ``everyone does it this way.''

This answers both the red pen's legitimacy and its limits. It is legitimate because agreement is real: the standard language used by more than a billion people is precious common property worth maintaining. It oversteps because there is more than one agreement. Cantonese has its conventions; Shanghainese has its own. Both are solidly established systems, not ``no agreement.'' The exam condemning Ah Chan should really say not ``this sentence is ungrammatical,'' but ``this sentence violates \textbf{the convention used in this exam}.''

Read another primary text, of a different kind. An old children's folk rhyme circulates across China. Its wording varies, but its skeleton stays the same:

\begin{quote}
Once there was a mountain. On the mountain stood a temple. In the temple an old monk was telling a story. What was the story? Once there was a mountain. On the mountain stood a temple\ldots{}
\end{quote}
Children giggle because they discover a mechanism: the story contains itself. Unless the storyteller stops, it can continue forever without ending.

Do not underestimate this rhyme. It points to the chapter's deepest discovery: grammar lets a structure \textbf{fit inside itself}. ``On the mountain stood a temple'' is a complete sentence, tucked whole inside ``the old monk tells a story,'' which is itself part of ``once there was a mountain.'' Sentences nest like dolls or opposing mirrors. Your mental grammar already knows this trick without anyone teaching it. Next we will see the enormous scientific dispute it has sparked.

\section{Is Grammar Assembled or Grown?}
We now have enough facts to compare genuinely different explanations of the same phenomena. First state what needs explaining, the evidential level with little to dispute: children everywhere, in Tokyo, Cairo, or Lima, master their native language's basic grammar around age three or four; nobody formally teaches them; they utter new sentences they have never heard; and their order and pace of acquisition are remarkably similar. The explanatory question is: \textbf{how does this precise apparatus get into every child's brain?}

\textbf{Explanation one: it comes preinstalled (generative grammar).}

In the 1950s, American linguist Noam Chomsky proposed an idea that transformed the discipline. What children hear is too sparse and messy---parents make slips, trail off, and speak unclearly---to let them infer an entire grammar from scratch. The human brain must therefore come with a ``language device,'' containing a basic blueprint shared by all human languages, which he called ``universal grammar.'' After birth, children need only use the speech they hear to set the switches to ``Chinese'' or ``English.'' On this account, acquiring grammar is less like learning to ride a bicycle than growing teeth: it develops when the time comes, while the environment supplies nourishment.

One of this account's strongest pieces of evidence comes from Central America. In the 1980s, Nicaragua first brought deaf children together for schooling. Previously scattered, they had no signing environment at home and did not know one another. Something happened that adults neither taught nor could teach: successive cohorts created a sign language on playgrounds and in dormitories. Younger arrivals refined it further, developing systematic grammatical structures. Today it is called Nicaraguan Sign Language, a genuine language. No teachers, no textbooks, no ``immortal'' academicians voting: grammar grew from children's interaction. It is hard to find a cleaner experiment answering whether teachers invent grammar.

The account's limits need stating too. Scholars have argued for decades about what the ``preinstalled blueprint'' actually contains, with no universally accepted inventory. And ``innate'' is easily mistaken for ``environment does not matter,'' which is plainly wrong. Without linguistic surroundings the apparatus develops nothing; the Nicaraguan children created language precisely through \textbf{interaction with one another}.

\textbf{Explanation two: it is inferred from use (usage-based theory).}

Other researchers, including teams led by psychologist Michael Tomasello, see no need for a preinstalled blueprint. They point to children's formidable general learning abilities: statistics, analogy, pattern-finding. From the vast speech they hear, children pan for frequent patterns like gold. ``Noun'' is a category extracted this way; ``word order'' is an extracted regularity. Grammar is the stock deposited by billions of uses. Their evidence: children master frequent words and constructions early and make more mistakes with rare irregularities. If the blueprint is preinstalled, why does learning speed track exposure so closely? On this account, learning grammar really is like learning to ride a bicycle; children are simply astonishingly quick.

This account also leaves something out of reach. It explains learning grammar from an existing language well, but Nicaragua's children had \textbf{no} existing language in front of them. Induction needs raw material; where was theirs? If it consisted of the ``urge to express something'' and one another's gestures, the leap from gestures to a language with systematic grammar still seems to require some distinctively human intuition for structure.

Both explanations hold part of the truth; neither reaches all of it. Remember their complete agreement on one point, precisely our title: whether grammar unfolds a preinstalled blueprint or settles from a flood of usage, \textbf{teachers did not invent it}. Teachers prune branches and erect signposts. The tree's seed lies either in genes or in each act of speaking, certainly not in the red pen.

One further misconception should be cleared away. Does this mean prescriptive grammar is entirely unnecessary? That too is easy to misjudge. Recall our three glues: order, inflection, function words. Languages everywhere use them, but no natural law requires negation to use ``two words'' or ``one.'' Agreements therefore really do need maintenance, or they drift. French ne\ldots{}pas has remained stable for centuries precisely because dictionaries and schools act as anchors. The real relationship between description and prescription is this: a grammar book is a map. Even the finest map is not the territory, but when you travel far, it is wise to bring one.

\section{Further Challenge: Infinitely Many Sentences}
Now we arrive at the chapter's weightiest theoretical question, hidden in the ``once there was a mountain'' rhyme.

Begin with a simple fact. Your vocabulary is finite: a few thousand words, perhaps ten thousand at most. Your possible lifetime sentences are infinite. In fact, almost any somewhat lengthy sentence you say now may be the first of its kind in human history: ``Yesterday I saw an orange cat perched on a parcel locker contemplating life.'' No grammar book has included it, yet it is perfectly legitimate. How can a few thousand finite parts assemble infinitely many sentences?

Early in the nineteenth century, German scholar Wilhelm von Humboldt expressed a famous idea: language makes \textbf{infinite use} of \textbf{finite means}. The central component of this ``infinity engine'' is the rhyme's mechanism: \textbf{recursion}. A structure fits inside another of the same kind; the resulting whole is still that kind of structure, so it can be inserted again and again, without grammatical limit.

See it work in sentences. ``I know'' is complete. Put it whole inside ``You know\ldots{}'': ``You know I know.'' Still complete. Insert again: ``He says you know I know.'' In theory, ``My mother says the teacher told him a classmate said the principal says you know I know'' can continue indefinitely. Grammar never calls a halt. Memory and patience do, not the rules.

What is extraordinary about recursion is that it is \textbf{learned once and usable for life}. You need no separate rule for each nested doll. Three- and four-year-olds already use it: ``This is the mouse chased by the dog that bit the cat.'' Without instruction, children acquire the ability to put sentences inside sentences. This is also where ``grammar is memorized'' breaks down entirely. Sentences cannot all be memorized; what can be learned is a \textbf{sentence-making machine}.

But watch an easily misjudged counterexample. If recursion has no grammatical limit, is every long sentence that can be assembled ``grammatical''? Consider: ``The dog that the cat bit ran away.'' Fine; with a little mental detour, you understand. Now: ``The mouse that the cat that the dog bit scratched ran away.'' Most people crash on the second loop, although its structure matches the earlier one and it is grammatically valid. What does this show? That ``grammatical'' and ``understandable'' differ. Grammar governs structural legitimacy; memory governs whether your brain can keep up. A sentence you cannot understand may overload working memory without violating any rule. Conversely, smooth, fluent speech may violate rules everywhere: listen to a recording of your own conversation. This distinction between structural legitimacy and processing feasibility is one of linguists' most important wrenches for opening the black box of ``language intuition.''

Go one step deeper: recursion may be among the most important boundaries between human language and animal communication. Bee dances and monkey alarm calls are sophisticated signaling systems, but \textbf{closed inventories}: one call corresponds to one situation, then ends. They cannot nest or build sentences. Every human child, however, uses an open-ended apparatus with unlimited output. Notice the wording: this is currently \textbf{one mainstream scientific view}, not a final verdict. Research continues into animals' rudimentary combinatorial capacities; crows and parrots occasionally surprise us. This is the stance we keep practicing: separate ``evidence shows,'' ``theory claims,'' and ``not yet settled.''

Return to Xunzi. He said names arise from ``agreement becoming custom.'' Correct, but only half the answer. Convention explains why rules take \textbf{this form} rather than another. It does not explain why every human society \textbf{has} rules, all containing recursion's engine. A collective vote cannot produce a heart's structure, nor ``sentences can contain sentences.'' That is something hundreds of thousands of years of evolution engraved in our species. Convention chooses the engine's fuel---gasoline, diesel, electricity, varying by place. But ``having an engine'' is humanity's shared factory equipment.

\section{Today: Chat Grammar and Talking Machines}
This ancient question plays out in your pocket every day.

Internet language is a high-speed accelerator of ``agreement becoming custom.'' Terms like yyds, shorthand for ``eternally the greatest,'' and juéjuézi, an emphatic ``absolutely amazing,'' went from nonexistent to shared understanding within weeks. No academy voted; no dictionary included them. Hundreds of millions used them first, and dictionaries hurried to catch up. This is Xunzi's process accelerated from centuries to weeks. Notice too how quickly new usage grows its own grammar. At first juéjuézi could appear almost anywhere; soon people tacitly learned its position, what it modified, and when its tone was ironic. Rules settle automatically through use, just as in older generations' dialects. You may dislike these words, but ``I dislike it'' and ``it has no grammar'' are different claims.

The ``faulty sentence'' question still sits on your exam paper. Now you can locate it more precisely. It tests not ``the laws of language itself,'' but ``the conventions of the particular community using standard written language.'' Those conventions deserve serious study: they are the public road for future résumés, contracts, and emails to strangers. But you can now mentally add a small line to the red pen's comment: this violates one set of conventions, not necessarily every system, and certainly does not mean its speaker ``cannot speak.'' A Guangdong child fluent in Cantonese and a child speaking only Mandarin each possess a complete, precise sentence-making machine. Only the settings differ.

Then comes the newest and most unsettling variable: machines have begun to speak. Today's large language models, including conversational assistants you can use, learn statistically from enormous amounts of text. Nobody gives them a grammar lesson or preinstalls ``universal grammar'' in their heads. They extract nearly impeccable grammar from the textual flood. This looks like a gigantic usage-based experiment. If induction really can extract grammar from use, has a machine using humanity's writings as raw material and a supercomputer as its pan found the gold? To an extent, yes. But does it \textbf{really understand}? Is there commitment behind its ``I promise''? Is its recursion understanding or imitation? Scientists and philosophers continue arguing fiercely. There is no standard answer here, only increasingly interesting questions.

One quiet corner also deserves attention. Across China, older people are slowly departing with their dialects; speaker numbers fall every year. Linguists rush to record them as biologists try to save endangered species, not from nostalgia but because each dialect is a unique underground palace. Cantonese's intact entering tones, Southern Min's layered literary and colloquial readings, Wu's soft voiced sounds preserve more than a thousand years of linguistic history. Once lost, they cannot be fully restored. Mandarin is a bridge; dialect is home. Bridges bring everyone together; home remembers where each person came from.

\section{Over to You: Who Should Hold the Red Pen?}
Return to the classroom with its ceiling fan and cicadas.

Ah Chan asks: ``Everyone where I'm from says `you leave first' that way. Are they all wrong?'' Teacher Zhou replies: ``In exams, follow the standard.''

Now answer a larger question: \textbf{if enough people are ``wrong'' for long enough, does wrong become right?}

Before concluding, weigh both sides again.

One says: language history is full of such overturned verdicts. In the Zuo Commentary's account of Duke Zhuang's tenth year, Cao Gui explains battle morale: ``At the first drumbeat, spirits rise; at the second, they weaken; at the third, they are exhausted.'' The word zài there means ``the second time''; today it generally means ``again.'' The old usage might now be marked wrong. English you was originally plural and respectful; singular thou was used less and less until it disappeared completely. Each era's ``incorrect usage'' may become the next era's standard. If established convention is grammar's sole legislature, every sentence the red pen condemns is only \textbf{the present} vote count, not an eternal judgment.

The other says: someone must count votes and repair roads. If every mouth is a legislature, what happens to exams and statutes? How does someone in Heilongjiang telephone someone in Guangxi? Abandoning standards does not liberate language; it dismantles a billion people's public road.

A deeper question lies underneath. Who holds the red pen today: increasingly authoritative dictionaries, or the ever-faster internet? When usage spreads nationwide in three days and is obsolete in three months, how long must ``custom'' take before it counts as established? If effective power over grammar passes to trending topics and algorithms, is that more democratic, or more dangerous?

There is no standard answer. But next time a red pen circles your writing, you can at least do something new: pick up the pen's logic and ask which set of conventions issued the judgment. Does it mean ``incoherent,'' or ``inappropriate here and now''? And if I have grounds to disagree, at what level can my reasons stand?

When you can ask that, the red pen remains in the teacher's hand, but grammar has returned to yours.
